\documentclass{article}
\usepackage[utf8]{inputenc}
\usepackage{indentfirst}
% \usepackage[letter, portrait, margin=1in]{geometry}
\usepackage[portrait, margin = 1in]{geometry}
\usepackage{array}
\usepackage{graphicx}
\usepackage{amsmath}
\usepackage{caption}
\usepackage{subcaption}
\usepackage{scicite}
\usepackage{textcomp}
\usepackage{siunitx}
\usepackage{hyperref}
\usepackage{wasysym}
\usepackage{setspace}
\hypersetup{
    colorlinks=true,
    linkcolor=cyan,
    filecolor=magenta,      
    urlcolor=cyan,
    pdftitle={AGU Abstract 2024}
    }
\graphicspath{ {/home/brynn/Code/OPT2k.jl/plots} }


\title{Arctic amplification of Little Ice Age cooling inferred from inversion of high-resolution benthic foraminiferal records}
\author{Brynnydd Hamilton, Geoffrey Gebbie, Wanyi Lu, Delia Oppo}
\date{July 2024}

\doublespacing
\begin{document}
\maketitle
Coupled climate models have suggested the presence of an Arctic amplification of Little Ice Age cooling, but it has not been conclusively shown to exist in paleoclimate reconstructions. The sparse and noisy nature of surface records (e.g. planktonic foraminifera, coral, or bivalve records) has complicated efforts to reconstruct variables of interest such as sea surface temperature. Benthic foraminifera record climate signals with higher fidelity than equivalent planktonic records, as they are not subject to high-frequency near-surface variability. In addition, they integrate the characteristics of their source water regions, providing more information about broader spatial scales. 
Here, we present the first North Atlantic surface history of Common Era sea surface temperatures and $\delta^{18}\mathrm{O}_{\mathrm{seawater}}$ inferred from benthic foraminiferal records via a global circulation inverse model. 
The methodology presented here accounts for age model uncertainty, as well as the lag time associated with the transit of waters from the surface to the sediment core sites.
Our reconstruction is within error of instrumental records, and reproduces a long-term cooling trend over the last millennium, in agreement with previous reconstructions based on terrestrial and marine records. In addition, a greater magnitude cooling is reconstructed in the Nordic Seas than in the subpolar gyre, indicative of an Arctic amplification of Little Ice Age cooling, previously suggested to exist in climate model simulations of the last millennium.

\end{document} 
%%% Local Variables:
%%% mode: latex
%%% TeX-master: t
%%% End:
